% GENERATED FILE. Edit canonical lessons, then run npm run generate:books.
% canonical-source-hash: fnv1a64:30d6d4b1f6b6e61a
% canonical-lessons: LA-C376-opus, LA-C376-opifex, LA-C376-artifex, LA-C376-operarius, LA-C376-operor

\chapter{Work and Workers}
\label{ch:la-work-and-workers}
\hlchaptermodality{376}

\begin{chapteropening}
\textbf{By the end of this chapter:} \emph{I can talk about a piece of work and say that I need something, and name a craftsman, an artist and a labourer at work.}

Complete the last lesson of chapter 376: I can talk about a piece of work and say that I need something, and name a craftsman, an artist and a labourer at work.
\end{chapteropening}

\section[opus, operis]{opus, operis — a work, a task; need (opus est)}

\label{lesson:LA-C376-opus}



\begin{quote}
Before the new one: say the Latin for unexpected, then the Latin for a movement.
\end{quote}

\subsection*{You'll want to know: opus, operis}
\textbf{opus, operis} — \textquotedblleft{}a work, a task; need (opus est)\textquotedblright{}.

\begin{quote}
Mihi pecūniā opus est.
\end{quote}

\begin{quote}
\emph{I need money.}
\end{quote}

\emph{opus, operis} n. \emph{Opus est} + dative of the person and ablative of the thing: \emph{mihi pecūniā opus est}, I need money. English \emph{operate} comes from its family, and \emph{opera} comes through Italian from the related \emph{opera}, effort.

\begin{cousinweb}
\emph{opus} is inherited from a Proto-Indo-European noun built on a root meaning \textquotedblleft{}to work, to produce\textquotedblright{} (EtymDB), with \emph{ops} \textquotedblleft{}power, means\textquotedblright{} as a cousin (EtymDB; Lewis \& Short). English \emph{operate} comes from its family, and \emph{opera} from the related \emph{opera} \textquotedblleft{}effort\textquotedblright{}.
\end{cousinweb}

\begin{grammarlens}[title={What you've built}]
One more word for telling, explaining and describing at B1.
\end{grammarlens}

\subsection*{Your turn}
\begin{itemize}
\raggedright
  \item \emph{Say it:} \emph{opus, operis}
  \item \emph{Say it:} \emph{opus, operis}, once more
  \item \emph{Say it:} a sentence of your own with \emph{opus, operis}
  \item \emph{Recall:} say the Latin for unexpected, then the Latin for a movement, then say \emph{opus, operis} again
\end{itemize}

\begin{tcolorbox}[breakable,colback=teal!4,colframe=teal!35!black,title={Before you move on}]
What does opus, operis mean? (\textquotedblleft{}A work, a task; need (opus est)\textquotedblright{}.) Say it once more.
\end{tcolorbox}
\section[opifex, opificis]{opifex, opificis — a craftsman, a worker}

\label{lesson:LA-C376-opifex}



\begin{quote}
Before the new one: say the Latin for a movement, then the Latin for a work.
\end{quote}

\subsection*{You'll want to know: opifex, opificis}
\textbf{opifex, opificis} — \textquotedblleft{}a craftsman, a worker\textquotedblright{}.

\begin{quote}
Opifex mēnsam novam fēcit.
\end{quote}

\begin{quote}
\emph{The craftsman made a new table.}
\end{quote}

\emph{opifex, opificis} m. and f. Built from \emph{opus}, work, and \emph{faciō}, to make: a work-maker.

\begin{cousinweb}
\emph{opifex} is built from \emph{opus} \textquotedblleft{}work\textquotedblright{} and \emph{faciō} \textquotedblleft{}to make\textquotedblright{} (Lewis \& Short; EtymDB puts the related \emph{ops} \textquotedblleft{}power, means\textquotedblright{} in place of \emph{opus}), with the ending \emph{-fex} of a maker.
\end{cousinweb}

\begin{grammarlens}[title={What you've built}]
One more word for telling, explaining and describing at B1.
\end{grammarlens}

\subsection*{Your turn}
\begin{itemize}
\raggedright
  \item \emph{Say it:} \emph{opifex, opificis}
  \item \emph{Say it:} \emph{opifex, opificis}, once more
  \item \emph{Say it:} a sentence of your own with \emph{opifex, opificis}
  \item \emph{Recall:} say the Latin for a movement, then the Latin for a work, then say \emph{opifex, opificis} again
\end{itemize}

\begin{tcolorbox}[breakable,colback=teal!4,colframe=teal!35!black,title={Before you move on}]
What does opifex, opificis mean? (\textquotedblleft{}A craftsman, a worker\textquotedblright{}.) Say it once more.
\end{tcolorbox}
\section[artifex, artificis]{artifex, artificis — an artist, a skilled craftsman; a master (of)}

\label{lesson:LA-C376-artifex}



\begin{quote}
Before the new one: say the Latin for a work, then the Latin for a craftsman.
\end{quote}

\subsection*{You'll want to know: artifex, artificis}
\textbf{artifex, artificis} — \textquotedblleft{}an artist, a skilled craftsman; a master (of)\textquotedblright{}.

\begin{quote}
Artifex statuam pulchram fēcit.
\end{quote}

\begin{quote}
\emph{The artist made a beautiful statue.}
\end{quote}

\emph{artifex, artificis} m. Built from \emph{ars} and \emph{faciō}: a maker by skill. \emph{Artifex dīcendī}, a master of speaking. English \emph{artifice} and \emph{artificial}.

\begin{cousinweb}
\emph{artifex} is built from \emph{ars, artis} \textquotedblleft{}skill\textquotedblright{} and \emph{faciō} \textquotedblleft{}to make\textquotedblright{} (EtymDB; Lewis \& Short). English \emph{artifice} and \emph{artificial} come from its family.
\end{cousinweb}

\begin{grammarlens}[title={What you've built}]
One more word for telling, explaining and describing at B1.
\end{grammarlens}

\subsection*{Your turn}
\begin{itemize}
\raggedright
  \item \emph{Say it:} \emph{artifex, artificis}
  \item \emph{Say it:} \emph{artifex, artificis}, once more
  \item \emph{Say it:} a sentence of your own with \emph{artifex, artificis}
  \item \emph{Recall:} say the Latin for a work, then the Latin for a craftsman, then say \emph{artifex, artificis} again
\end{itemize}

\begin{tcolorbox}[breakable,colback=teal!4,colframe=teal!35!black,title={Before you move on}]
What does artifex, artificis mean? (\textquotedblleft{}An artist, a skilled craftsman; a master (of)\textquotedblright{}.) Say it once more.
\end{tcolorbox}
\section[operārius, operāriī]{operārius, operāriī — a labourer, a workman}

\label{lesson:LA-C376-operarius}



\begin{quote}
Before the new one: say the Latin for a craftsman, then the Latin for an artist.
\end{quote}

\subsection*{You'll want to know: operārius, operāriī}
\textbf{operārius, operāriī} — \textquotedblleft{}a labourer, a workman\textquotedblright{}.

\begin{quote}
Operāriī tōtum diem in agrō labōrant.
\end{quote}

\begin{quote}
\emph{The labourers work all day in the field.}
\end{quote}

\emph{operārius, operāriī} m., a man who works with his hands for a wage, built on \emph{opera}, work, effort.

\begin{cousinweb}
\emph{operārius} is built on \emph{opera} \textquotedblleft{}work, effort\textquotedblright{} (Lewis \& Short), a word of the same family as \emph{opus, operis} \textquotedblleft{}work\textquotedblright{}, with the ending \emph{-ārius} of a person busy with something.
\end{cousinweb}

\begin{grammarlens}[title={What you've built}]
One more word for telling, explaining and describing at B1.
\end{grammarlens}

\subsection*{Your turn}
\begin{itemize}
\raggedright
  \item \emph{Say it:} \emph{operārius, operāriī}
  \item \emph{Say it:} \emph{operārius, operāriī}, once more
  \item \emph{Say it:} a sentence of your own with \emph{operārius, operāriī}
  \item \emph{Recall:} say the Latin for a craftsman, then the Latin for an artist, then say \emph{operārius, operāriī} again
\end{itemize}

\begin{tcolorbox}[breakable,colback=teal!4,colframe=teal!35!black,title={Before you move on}]
What does operārius, operāriī mean? (\textquotedblleft{}A labourer, a workman\textquotedblright{}.) Say it once more.
\end{tcolorbox}
\section[operor, operārī]{operor, operārī — to work, to be busy (at)}

\label{lesson:LA-C376-operor}



\begin{quote}
Before the new one: say the Latin for an artist, then the Latin for a labourer.
\end{quote}

\subsection*{You'll want to know: operor, operārī}
\textbf{operor, operārī} — \textquotedblleft{}to work, to be busy (at)\textquotedblright{}.

\begin{quote}
Fabrī in officīnā operābantur.
\end{quote}

\begin{quote}
\emph{The smiths were at work in the workshop.}
\end{quote}

Deponent: \emph{operor, operārī, operātus sum}, often with the dative of the task. English \emph{operate}.

\begin{cousinweb}
\emph{operor} is built on \emph{opus, operis} \textquotedblleft{}work\textquotedblright{} (EtymDB; Lewis \& Short). English \emph{operate} comes from it.
\end{cousinweb}

\begin{grammarlens}[title={What you've built}]
One more word for telling, explaining and describing at B1.
\end{grammarlens}

\subsection*{Your turn}
\begin{itemize}
\raggedright
  \item \emph{Say it:} \emph{operor, operārī}
  \item \emph{Say it:} \emph{operor, operārī}, once more
  \item \emph{Say it:} a sentence of your own with \emph{operor, operārī}
  \item \emph{Recall:} say the Latin for an artist, then the Latin for a labourer, then say \emph{operor, operārī} again
\end{itemize}

\begin{tcolorbox}[breakable,colback=teal!4,colframe=teal!35!black,title={Before you move on}]
What does operor, operārī mean? (\textquotedblleft{}To work, to be busy (at)\textquotedblright{}.) Say it once more.
\end{tcolorbox}
